\bookchapter{Capítulo 0: Hola, antes de empezar}

Antes de la exploración formal, hablemos de algunas reglas básicas para asegurarnos de que entiendes lo que viene. No te apresures, afilar el hacha no retrasa el trabajo de cortar leña.

\booksection{¿Qué puede hacer este libro por ti? (Y qué no)}

Puede ayudarte a entender sistemáticamente cómo usar tu móvil, a comprender la lógica detrás de las funciones, y ya seas del equipo Android o del equipo Apple, la mayoría de las ideas son aplicables. Te guiará en ejercicios prácticos para que pases de “saber usar” a “saber usar bien”.

No es el manual de un modelo específico de móvil, ni son noticias de tecnología sobre los últimos productos, ni es una guía para que te enganches al móvil. Sinceramente, si lo que esperas es “cómo ver videos en tu móvil más fluidamente”, puede que te decepcione — nos importa más cómo hacer que el móvil realmente te sirva, no lo contrario.

\booksection{Explicación de nuestras “claves”}

Para explicarlo más claramente, el libro usa algunas marcas pequeñas, estas marcas son como códigos tácitos entre nosotros:

\begin{itemize}
  \item \textbf{[Nombre del botón]}: representa un botón específico en la pantalla o el nombre de una App, por ejemplo [Configuración], [Aceptar].
  \item \textbf{>}: representa el orden de operación, por ejemplo "tocar [Configuración] > [Sonido]" significa primero tocar Configuración, luego Sonido.
  \item \textbf{Negrita}: ¡muy importante! Puede ser un concepto clave o una advertencia de seguridad, por favor presta atención.
  \item \emph{Cursiva}: algunas ideas complementarias o metáforas para ayudarte a entender mejor.
\end{itemize}

\booksection{Android y iOS, hablamos de ambos}

Adivina qué, algunas cosas tienen nombres diferentes pero hacen lo mismo. Por ejemplo: aplicación (Android) y App (iOS), Configuración (ambos se llaman así), Gestor de archivos (Android) y Archivos (iOS). Saber que son lo mismo es suficiente, no te confundas con los términos.

Además, ahora en el campo de Android también está el sistema HarmonyOS de Huawei, que tiene sus propias carácterísticas, lo trataremos especialmente más adelante. En resumen, sin importar qué móvil uses, este libro puede ayudarte.

\booksection{Capacidad del móvil: ¿por qué compré 128G pero solo hay menos de 110G disponibles?}

Las unidades de capacidad las conoces: 1TB = 1024GB, 1GB = 1024MB, 1MB = 1024KB... pero aquí hay una trampa: los fabricantes de discos suelen calcular en base 1000, mientras que el sistema calcula en base 1024, así que el disco “128G” que compras, la capacidad real “se reducirá”.

Pero las razones principales por las que se reduce la capacidad de tu móvil son tres:

\begin{enumerate}
  \item \textbf{El sistema mismo ocupa espacio}: el sistema operativo es como el muro de carga de una casa, debe existir. Los sistemas actuales fácilmente ocupan 10-20GB, esto es normal.
  \item \textbf{Aplicaciones preinstaladas del fabricante}: las Apps que vienen con el móvil nuevo también ocupan espacio, algunas se pueden desinstalar, otras no.
  \item \textbf{Pérdidas por partición y formateo}: por razones técnicas, como el forro de una mochila también ocupa volumen.
\end{enumerate}

\begin{tcbinfo}{Pequeño consejo para comprar móvil}
Ahora los móviles de 128GB básicamente no son suficientes, 256GB para empezar es más prudente. Si te gusta tomar fotos, jugar, se recomienda directamente 512GB o más. Recuerda, el espacio de almacenamiento, mejor que sobre y no que falte.

\end{tcbinfo}

\begin{tcbpractice}{Manos a la obra: mira la "identidad" de tu móvil}
Abre [Configuración], busca Acerca del teléfono (o una opción similar), anota: modelo del móvil, versión del sistema, espacio de almacenamiento total y espacio libre. Este es el primer paso para conocerlo. Aprovecha para revisar la salud de la batería, para tener una idea.

\end{tcbpractice}

\booksection{Tu móvil es una “biblioteca de activos digitales”}

Antes de empezar con el conocimiento específico, primero establezcamos una cognición central:\textbf{el móvil hace tiempo que no es solo una herramienta de comunicación, es nuestra biblioteca de activos digitales más importante}. Contactos, registros de chat, fotos familiares, documentos de trabajo, Apps bancarias, códigos de pago, datos de salud, documentos electrónicos... todos estos datos juntos, su valor a menudo supera al del propio móvil.

\begin{tcbinfo}{¿Por qué es importante?}
Si pierdes o se rompe el móvil, no solo pierdes hardware de varios miles,\textbf{sino que pueden ser años de datos irregenerables acumulados}—fotos del crecimiento de los hijos, audios de familiares fallecidos, contratos importantes sin respaldo. Los hábitos científicos de organización y protección pueden protegerte en tres dimensiones:

\begin{itemize}
  \item \textbf{Mejora la eficiencia}: una clasificación de archivos estándarizada hace que el sistema busque más rápido, el móvil no se atasca a largo plazo.
  \item \textbf{Prolonga la vida útil}: una gestión razonable del almacenamiento reduce la carga de lectura/escritura del chip, la batería y el almacenamiento duran más.
  \item \textbf{Garantiza la seguridad de los datos}: gestión por capas + copia de seguridad periódica + reducción de permisos, mantiene las "catástrofes de datos" alejadas.
\end{itemize}

\end{tcbinfo}

\booksubsection{Nuestra filosofía central: “Prevención ante todo, combinación de prevención y tratamiento”}

Como cuidar la salud no puede esperar a enfermarse para ver al médico,\textbf{la gestión del móvil también debe empezar por los “hábitos diarios”}. Lo que este libro enfatiza repetidamente no es una técnica avanzada concreta, sino estos principios sencillos:

\begin{itemize}
  \item \textbf{Prevención ante todo}: instalar Apps solo desde tiendas oficiales, desensibilizar fotos antes de enviarlas, no repetir contraseñas—estos "pequeños gestos" tienen un coste bajísimo, pero bloquean la mayoría de los riesgos.
  \item \textbf{Revisión periódica}: organizar la carpeta de descargas una vez al mes, auditoría de permisos cada trimestre, copia de seguridad completa cada semestre—un ritmo fijo es cien veces más efectivo que apagar fuegos.
  \item \textbf{Dejar una salida}: activar la sincronización en la nube, anotar bien las cuentas y contraseñas, escribir una "lista de legado digital"—para asegurar que cuando ocurra algo imprevisto, tú y tu familia puedan manejarlo con calma.
\end{itemize}
